\documentclass[10pt]{article}
\usepackage[a4paper,margin=22mm]{geometry}
\usepackage[T1]{fontenc}
\usepackage{lmodern}
\usepackage[hidelinks]{hyperref}
\setlength{\parindent}{0pt}
\setlength{\parskip}{5pt}
\setlength{\emergencystretch}{2em}
\pagestyle{empty}

\begin{document}
\raggedright

\begin{flushright}
Department of Computer Science and Engineering\\
SAGE University\\
Indore 452020, Madhya Pradesh, India
\end{flushright}

The Editors\\
\emph{Journal of Seismology}\\
Springer Nature

\medskip
Dear Editors,

We submit for your consideration our manuscript, ``A versioned,
leakage-audited USGS \emph{Did You Feel It?} benchmark for machine-learning
evaluation of severe felt intensity.'' We request handling as a data or software
paper documenting a public seismological resource; if that label is unavailable
in Snapp, we request consideration as an Original Paper.

Machine-learning studies of macroseismic felt intensity commonly use bespoke
datasets and random record-level splits. Their results can be difficult to
reproduce and vulnerable to leakage between related earthquake events. Our
manuscript addresses this problem by releasing a frozen, single-source USGS
DYFI benchmark, a six-category leakage audit, outcome-independent
sequence-grouped roles, fixed baselines, calibration-aware scoring, and a
protected temporal evaluation.

The frozen \(M\geq5\) snapshot contains 7,520 source records, 2,362 eligible
events, and a 430-event temporal holdout. The observed random-versus-sequence
split differences are small and reported descriptively; no winning model,
operational forecasting benefit, causal effect, or cross-system generalization
is claimed. All deviations and unexecuted branches are disclosed.

The work fits the \emph{Journal of Seismology} as an observational
macroseismology and citizen-science data/software resource. It directly aligns
with the journal's stated interest in machine learning, seismological datasets,
software tools, and publicly available resources supporting transparent,
reproducible earthquake science.

Data, code, frozen results, and reproduction instructions are available at
\url{https://github.com/akssha74/dyfi-benchmark-reliability}. A clean clone
reproduces the 15-page manuscript, passes all 106 manuscript checks and 110
benchmark tests, and reconstructs the recorded results and holdout predictions.

The manuscript is original, is not under consideration elsewhere, and has not
been published previously. The authors declare no competing interests. The
study uses aggregate public-domain event metadata, involves no human
participants, and redistributes no respondent-level identifying information.
Generative-AI assistance is disclosed in Methods and all assisted work was
verified by the authors.

\medskip
Sincerely,

Akshay Sharma (corresponding author) and Lalji Prasad\\
Department of Computer Science and Engineering\\
SAGE University, Indore, India\\
\href{mailto:akssha74@gmail.com}{akssha74@gmail.com}

\end{document}
